\chapter{PENDAHULUAN}

\section{Latar Belakang}
Bab Pendahuluan adalah bab terpendek dalam sebuah tesis dan merupakan bab yang paling penting. Satu-satunya tujuan penulisan bab ini adalah untuk memperkenalkan kepada pembaca tentang apa yang akan dipaparkan dalam tesis \citep{russell2020}.

Jelaskan lebih dahulu apa yang memotivasi penelitian yang akan dilakukan dan mengapa penelitian pada domain \textit{[nama domain aplikasi]} tersebut perlu dilakukan. Paparkan secara sistematis gejala dan peristiwa yang menimbulkan permasalahan untuk diteliti, serta sebutkan penelitian-penelitian terdahulu yang relevan pada bidang yang akan diteliti \citep{goh2009}.

\subsection{Ruang lingkup permasalahan yang diteliti}
Pada bagian ini penulis menjelaskan secara ringkas gambaran umum permasalahan penelitian, batasan awal ruang lingkupnya, serta perbedaan kontribusi penelitian ini dibandingkan dengan penelitian-penelitian sebelumnya yang telah disebutkan.

\section{Rumusan Masalah}
Rumusan masalah dapat ditulis dalam bentuk \textit{research question} atau \textit{problem statement}. Contoh rumusan masalah berbentuk pertanyaan penelitian:
\begin{enumerate}
    \item Bagaimana pengaruh \textit{[variabel/metode]} terhadap kinerja \textit{[tugas/permasalahan]} pada \textit{[domain aplikasi]}?
    \item Bagaimana perbandingan kinerja pendekatan yang diusulkan terhadap metode \textit{baseline} yang ada?
\end{enumerate}

\section{Tujuan Penelitian}
Tujuan penelitian ini adalah untuk merancang, mengimplementasikan, dan mengevaluasi \textit{[nama metode/algoritma]} guna menyelesaikan permasalahan yang telah dirumuskan pada Subbab 1.2.

\section{Manfaat Penelitian}
Manfaat yang diharapkan dari penelitian ini antara lain memberikan kontribusi teoritis pada bidang \textit{[bidang keilmuan]} serta manfaat praktis berupa \textit{[manfaat terapan, misalnya: purwarupa sistem, rekomendasi kebijakan, dsb.]}.

\section{Batasan Masalah}
Agar penelitian tetap fokus, ruang lingkup penelitian dibatasi pada hal-hal berikut:
\begin{enumerate}
    \item Penelitian dibatasi pada disiplin ilmu data science dan kecerdasan buatan.
    \item Data yang digunakan adalah \textit{[nama/sumber dataset]}.
    \item Evaluasi dilakukan menggunakan metrik \textit{[nama metrik]} dan tidak mencakup \textit{[hal yang dikecualikan]}.
\end{enumerate}

\section{Sistematika Penulisan}
Sistematika penulisan tesis ini terdiri atas lima bab. Bab 1 berisi pendahuluan. Bab 2 membahas tinjauan pustaka dan landasan teori. Bab 3 menguraikan metodologi penelitian. Bab 4 menyajikan hasil dan pembahasan. Bab 5 berisi kesimpulan dan saran.
